% ═══════════════════════════════════════════════════════════════════════════
%  فصل ۶ — ایران عصر صفوی تا قاجار: مرزنشینان و همسایگان
% ═══════════════════════════════════════════════════════════════════════════
\chapter{ایران عصر صفوی تا قاجار: مرزنشینان و همسایگان}
\label{ch:early-modern-iran}

\begin{tcolorbox}[mybox, title={چکیده‌ی فصل}]
این فصل به بررسی نجات‌طلبی و مداخلات خارجی در ایران از آغاز دوره‌ی صفوی (۹۰۷ ق / ۱۵۰۱ م) تا پایان دوره‌ی قاجار (۱۳۰۴ ش / ۱۹۲۵ م) می‌پردازد. تمرکز بر سه محور اصلی است: (۱) روابط پیچیده‌ی مرزنشینان قفقاز (به‌ویژه گرجی‌ها و ارمنی‌ها) با قدرت‌های ایرانی، عثمانی و روسی؛ (۲) نقش قبایل عرب خلیج‌فارس و شیخ خزعل؛ (۳) الگوی از دست‌رفتن سرزمین‌ها در پی مداخلات خارجی و جنگ‌های ایران و روسیه. از مدل شش‌وجهی و منابع آرشیوی فارسی، عثمانی، روسی و بریتانیایی بهره گرفته می‌شود.
\end{tcolorbox}

% ─────────────────────────────────────────────────────────────────────────
\section{دوره‌ی صفوی: تشیع، هویت ملی و خطوط شکاف نوین}
\label{sec:ch06-safavid}
% ─────────────────────────────────────────────────────────────────────────

\subsection{تأسیس صفویه و ایجاد هویت ایرانی-شیعی}
\label{subsec:ch06-safavid-identity}

تأسیس دولت صفوی توسط \histref{شاه اسماعیل اول}{۹۰۷ ق / ۱۵۰۱ م} نقطه‌ی عطفی در تاریخ ایران بود: برای نخستین بار پس از فتح عرب، یک دولت مستقل ایرانی با هویت مذهبی متمایز (تشیع دوازده‌امامی) تأسیس شد. این تحول، خطوط شکاف جدیدی ایجاد کرد:

\begin{itemize}[nosep, rightmargin=1em]
  \item \textbf{شکاف مذهبی جدید:} تشیع اجباری، اقلیت سنّی‌مذهب — به‌ویژه در کردستان، بلوچستان و مناطق ترکمن‌نشین — را ناراضی کرد.%
  \footnote{نوایی، عبدالحسین. \textit{شاه اسماعیل صفوی}. ج ۱. بنیاد فرهنگ ایران، ۱۳۴۷.}
  \item \textbf{شکاف قومی-مذهبی در قفقاز:} گرجی‌های مسیحی و ارمنی‌ها، با وجود حضور در ساختار قدرت صفوی، هم‌کیشان خود در اروپا و بیزانس/روسیه را «نزدیک‌تر» احساس می‌کردند.%
  \footnote{\lr{Floor, Willem, and Edmund Herzig, eds. \textit{Iran and the World in the Safavid Age}. I.B. Tauris, 2012, pp.~252--280.}}
  \item \textbf{رقابت صفوی-عثمانی:} عثمانی‌ها به‌عنوان رقیب سنّی‌مذهب، از هر شکاف داخلی ایران بهره‌برداری می‌کردند — و بالعکس.
\end{itemize}

\subsection{قفقاز: گرجی‌ها میان ایران، عثمانی و روسیه}
\label{subsec:ch06-georgia}

\subsubsection{گرجستان در نظام صفوی}

گرجستان در دوره‌ی صفوی به‌عنوان \concept{ایالت وابسته} در ساختار شاهنشاهی ادغام شده بود. شاهان صفوی سیاست پیچیده‌ای در قبال گرجستان داشتند:

\begin{enumerate}[nosep, label=\textbf{\arabic*.}]
  \item \textbf{جذب نخبگان:} شاهزادگان گرجی (مانند خاندان باگراتیونی) پس از تغییر مذهب به تشیع، در مناصب بالای دربار صفوی جای می‌گرفتند. نمونه‌ی بارز: \histref{الله‌وردی‌خان}{ت. ۱۰۲۲ ق / ۱۶۱۳ م} سپهسالار شاه‌عباس اول.%
  \footnote{\lr{Matthee, Rudi. \textit{Persia in Crisis: Safavid Decline and the Fall of Isfahan}. I.B. Tauris, 2012, pp.~35--42.}}
  \item \textbf{سرکوب شورش‌ها:} شاه‌عباس اول در سال‌های ۱۰۲۳–۱۰۲۵ ق / ۱۶۱۴–۱۶۱۶ م لشکرکشی‌های سنگینی به کاخِتی و کارتلی کرد و حدود ۶۰–۱۰۰ هزار گرجی را کوچ اجباری داد.%
  \footnote{\lr{Blow, David. \textit{Shah Abbas: The Ruthless King Who Became an Iranian Legend}. I.B. Tauris, 2009, pp.~178--195.}}
  \item \textbf{نظام «غلام‌سازی»:} کودکان گرجی و ارمنی اسیر شده، در دربار صفوی تربیت و به مناصب نظامی و اداری می‌رسیدند — نوعی \concept{مملوک‌سازی ایرانی}.
\end{enumerate}

\subsubsection{نجات‌طلبی گرجی‌ها: رو به اروپا، رو به روسیه}

\begin{tcolorbox}[casebox, title={نمونه‌ی تاریخی: پادشاه تئیمورز دوم و دعوت از روسیه}]
\histref{تئیمورز دوم}{ح. ۱۱۳۰–۱۱۷۶ ق / ۱۷۱۸–۱۷۶۲ م} پادشاه کاخِتی و سپس تمام گرجستان شرقی، نمونه‌ی بارز نجات‌طلبی است. او در شرایطی که گرجستان میان فشار عثمانی از غرب و ضعف صفویه/افشاریه از شرق قرار داشت، بارها از \concept{روسیه‌ی تزاری} درخواست حمایت نظامی کرد.

\begin{itemize}[nosep, rightmargin=1em]
  \item \textbf{۱۱۴۴ ق / ۱۷۳۱ م:} تئیمورز نامه‌هایی به دربار روسیه نوشت و خواستار اعزام نیروی نظامی برای حمایت از گرجستان در برابر عثمانی و ایران شد.%
  \footnote{\lr{Suny, Ronald Grigor. \textit{The Making of the Georgian Nation}. 2nd ed. Indiana UP, 1994, pp.~50--55.}}
  \item \textbf{۱۱۷۶ ق / ۱۷۶۲ م:} پسر تئیمورز، \histref{ایراکلی دوم}{ح. ۱۱۷۴–۱۲۱۲ ق / ۱۷۶۰–۱۷۹۸ م}، سیاست پدر را ادامه داد و در نهایت \concept{عهدنامه‌ی گئورگیفسک} (۱۱۹۷ ق / ۱۷۸۳ م) را با روسیه امضا کرد — عهدنامه‌ای که گرجستان را تحت‌الحمایه‌ی روسیه قرار داد.%
  \footnote{\lr{Lang, David Marshall. \textit{The Last Years of the Georgian Monarchy, 1658--1832}. Columbia UP, 1957, pp.~183--210.}}
\end{itemize}
\end{tcolorbox}

\begin{tcolorbox}[casebox, title={تحلیل شش‌وجهی: نجات‌طلبی گرجی‌ها (۱۷۸۳)}]
\begin{itemize}[nosep, rightmargin=1em]
  \item \textbf{ساختاری:} گرجستان در فشار دوگانه‌ی عثمانی (تهاجم و اسارت) و ایران (سلطه‌ی سیاسی و تغییر مذهب اجباری) قرار داشت. پس از سقوط صفویه (۱۱۳۵ ق / ۱۷۲۲ م) و هرج‌ومرج دوره‌ی افشاری-زندی، حمایت مرکزی از گرجستان ناچیز شد.
  \item \textbf{حقوقی-هنجاری:} عهدنامه‌ی گئورگیفسک ظاهراً توافقی دوجانبه بود، اما روسیه تعهداتش (حمایت نظامی) را عملاً زیر پا گذاشت — نمونه: عدم کمک هنگام حمله‌ی \enemy{آقامحمدخان قاجار} به تفلیس (۱۲۱۰ ق / ۱۷۹۵ م).%
  \footnote{\lr{Suny, 1994, pp.~59--63.}}
  \item \textbf{اقتصادی:} گرجستان منابع معدنی و کشاورزی ارزشمندی داشت و راه ترانزیت مهمی بود. روسیه از این منافع بهره برد.
  \item \textbf{نظامی:} ایراکلی دوم ارتش کوچکی داشت (حدود ۵–۱۰ هزار نفر) که توانایی مقابله با نه عثمانی و نه ایران را نداشت.
  \item \textbf{روان‌شناختی:} هویت مسیحی ارتدوکس گرجی‌ها، روسیه‌ی مسیحی را «ناجی طبیعی» و ایران و عثمانی مسلمان را «ستم‌گر ذاتی» جلوه می‌داد — نمونه‌ی بارز \concept{تفکر سیاه-سفید}.
  \item \textbf{ژئوپولیتیک:} روسیه به قفقاز به‌عنوان دروازه‌ی نفوذ به ایران و عثمانی نیاز داشت. «حمایت از مسیحیان» بهانه‌ای عالی بود.
  \item \textbf{نتیجه (طیف):} \textbf{D–E}. گرجستان در ۱۲۱۵ ق / ۱۸۰۱ م رسماً به روسیه الحاق شد — نه به‌عنوان «متحد» بلکه به‌عنوان «استان». خاندان باگراتیونی از قدرت برکنار شد و گرجی‌ها تا ۱۹۱۸ هیچ حاکمیت مستقلی نداشتند.%
  \footnote{\lr{Rayfield, Donald. \textit{Edge of Empires: A History of Georgia}. Reaktion Books, 2012, pp.~248--285.}}
\end{itemize}
\end{tcolorbox}

\subsection{ارمنی‌ها: «ملت وفادار» یا «پل ارتباطی با بیگانه»؟}
\label{subsec:ch06-armenians}

ارمنی‌های ایران نقش پیچیده‌ای در تاریخ نجات‌طلبی داشتند. در دوره‌ی صفوی، ارمنی‌های جلفای اصفهان از ثروت و آزادی نسبی مذهبی برخوردار بودند.%
\footnote{\lr{Herzig, Edmund. "The Armenian Merchants of New Julfa, Isfahan." In \textit{Safavid Persia}, ed. Charles Melville. I.B. Tauris, 1996, pp.~298--320.}}

با این حال، \thinker{ایسرائل اوری} (\lr{Israel Ori})، تاجر و دیپلمات ارمنی، در اواخر سده‌ی هفدهم سفرهایی به اروپا کرد و از پاپ، لویی چهاردهم و شاهزاده‌ی پفالتس درخواست لشکرکشی به ایران برای «نجات ارمنیان» کرد.%
\footnote{\lr{Bournoutian, George A. "Eastern Armenia from the Seventeenth Century to the Russian Annexation." In \textit{The Armenian People from Ancient to Modern Times}, ed. Hovannisian. Macmillan, 1997, vol.~2, pp.~81--107.}}

\begin{longtable}{
  >{\raggedleft\arraybackslash\bfseries}p{2.5cm}
  >{\raggedleft\arraybackslash}p{2.5cm}
  >{\raggedleft\arraybackslash}p{3cm}
  >{\raggedleft\arraybackslash}p{2.5cm}
  >{\raggedleft\arraybackslash}p{2.5cm}}
\caption{نجات‌طلبی ارمنی‌ها در دوره‌ی صفوی-قاجار}\label{tab:ch06-armenian}\\
\toprule
\rowcolor{PrimaryDeep}
\textcolor{white}{\textbf{فعال/رویداد}} &
\textcolor{white}{\textbf{تاریخ}} &
\textcolor{white}{\textbf{مخاطب}} &
\textcolor{white}{\textbf{خواسته}} &
\textcolor{white}{\textbf{نتیجه}} \\
\midrule
\endfirsthead
\toprule
\rowcolor{PrimaryDeep}
\textcolor{white}{\textbf{فعال/رویداد}} &
\textcolor{white}{\textbf{تاریخ}} &
\textcolor{white}{\textbf{مخاطب}} &
\textcolor{white}{\textbf{خواسته}} &
\textcolor{white}{\textbf{نتیجه}} \\
\midrule
\endhead
\bottomrule
\endlastfoot

ایسرائل اوری &
۱۶۹۹–۱۷۱۱ &
پاپ، فرانسه، آلمان &
لشکرکشی به ایران &
شکست (مرگ اوری)%
\footnote{\lr{Bournoutian, 1997, pp.~91--93.}} \\
\rowcolor{NeutralLight}
ملیک‌های قراباغ &
قرن ۱۸ &
پطر کبیر (روسیه) &
حمایت نظامی &
لشکرکشی ناکام پطر به قفقاز%
\footnote{\lr{Swietochowski, Tadeusz. \textit{Russia and Azerbaijan: A Borderland in Transition}. Columbia UP, 1995, pp.~4--12.}} \\
یوسف امین &
۱۷۶۰–۱۷۷۰ &
کاترین دوم (روسیه) &
استقلال ارمنستان &
عدم اقدام عملی \\
\rowcolor{NeutralLight}
نجات‌طلبی عمومی &
۱۸۰۴–۱۸۲۸ &
ارتش روسیه &
پیوستن به روسیه &
الحاق ارمنستان شرقی (ترکمنچای)%
\footnote{\lr{Bournoutian, George A. \textit{The Khanate of Erevan under Qajar Rule, 1795--1828}. Mazda, 1992.}} \\

\end{longtable}

% ─────────────────────────────────────────────────────────────────────────
\section{جنگ‌های ایران و روسیه: نجات‌طلبی و از دست‌رفتن قفقاز}
\label{sec:ch06-russo-iranian}
% ─────────────────────────────────────────────────────────────────────────

\subsection{جنگ اول (۱۲۱۸–۱۲۲۸ ق / ۱۸۰۴–۱۸۱۳ م)}
\label{subsec:ch06-first-war}

جنگ اول ایران و روسیه در بستر الحاق یک‌جانبه‌ی گرجستان شرقی توسط روسیه (۱۸۰۱) آغاز شد. نکته‌ی مهم آن است که \textbf{بخش بزرگی از جمعیت مسیحی قفقاز — گرجی‌ها و ارمنی‌ها — از ارتش روسیه استقبال کردند} و حتی در آن خدمت نمودند.%
\footnote{\lr{Atkin, Muriel. \textit{Russia and Iran, 1780--1828}. U of Minnesota P, 1980, pp.~94--130.}}

\begin{tcolorbox}[legalbox, title={عهدنامه‌ی گلستان (۱۲۲۸ ق / ۱۸۱۳ م): پیامدهای ارضی}]
بر اساس عهدنامه‌ی گلستان، ایران خانات زیر را به روسیه واگذار کرد:
\begin{itemize}[nosep, rightmargin=1em]
  \item باکو، شیروان، شکی، قره‌باغ، گنجه
  \item دربند، کوبا، تالش
  \item گرجستان شرقی (تأیید الحاق ۱۸۰۱)
  \item ابخازیا، مینگرلیا، ایمرتی
\end{itemize}
\vspace{4pt}
نکته‌ی حقوقی: عهدنامه‌ی گلستان تحت فشار نظامی و با میانجی‌گری بریتانیا امضا شد. ایران حق داشتن نیروی دریایی در دریای خزر را نیز از دست داد.%
\footnote{هدایت، رضاقلی‌خان. \textit{روضة الصفای ناصری}. ج ۹. اساطیر، ۱۳۸۰، ص ۴۵۵--۴۸۰.}
\end{tcolorbox}

\subsection{جنگ دوم (۱۲۴۱–۱۲۴۳ ق / ۱۸۲۶–۱۸۲۸ م) و عهدنامه‌ی ترکمنچای}
\label{subsec:ch06-second-war}

جنگ دوم با تحریک بریتانیا و فتوای جهاد علمای تبریز آغاز شد اما با شکست سنگین ایران و عهدنامه‌ی ترکمنچای (۱۲۴۳ ق / ۱۸۲۸ م) پایان یافت.%
\footnote{\lr{Amanat, Abbas. \textit{Pivot of the Universe: Nasir al-Din Shah and the Iranian Monarchy, 1831--1896}. I.B. Tauris, 1997, pp.~11--28.}}

\begin{figure}[htbp]
\centering
\begin{tikzpicture}[
  every node/.style={font=\scriptsize},
  territory/.style={
    draw=PrimaryMid, fill=BgFail!40, rounded corners=2pt,
    text width=2.5cm, align=center, minimum height=0.8cm
  },
  remaining/.style={
    draw=AccentGreen, fill=BgSuccess!40, rounded corners=2pt,
    text width=2.5cm, align=center, minimum height=0.8cm
  },
]

% عنوان
\node[font=\normalsize\bfseries, color=PrimaryDeep] at (6,7.5) {\rl{سرزمین‌های ازدست‌رفته‌ی ایران در قفقاز}};

% گلستان (۱۸۱۳)
\node[font=\small\bfseries, color=AccentRed] at (3,6.5) {\rl{گلستان ۱۸۱۳}};
\node[territory] at (1,5.5) {\rl{باکو}};
\node[territory] at (3,5.5) {\rl{شیروان}};
\node[territory] at (5,5.5) {\rl{قره‌باغ}};
\node[territory] at (1,4.3) {\rl{گنجه}};
\node[territory] at (3,4.3) {\rl{شکی}};
\node[territory] at (5,4.3) {\rl{دربند}};
\node[territory] at (1,3.1) {\rl{گرجستان شرقی}};
\node[territory] at (3,3.1) {\rl{کوبا}};
\node[territory] at (5,3.1) {\rl{تالش (بخشی)}};

% ترکمنچای (۱۸۲۸)
\node[font=\small\bfseries, color=AccentRed] at (9,6.5) {\rl{ترکمنچای ۱۸۲۸}};
\node[territory] at (8,5.5) {\rl{ایروان}};
\node[territory] at (10,5.5) {\rl{نخجوان}};
\node[territory] at (9,4.3) {\rl{اردوباد}};

% آنچه ماند
\node[font=\small\bfseries, color=AccentGreen] at (6,1.8) {\rl{باقی‌مانده}};
\node[remaining] at (4,0.8) {\rl{آذربایجان ایران}};
\node[remaining] at (8,0.8) {\rl{گیلان و مازندران}};

% فلش
\draw[->, >=Stealth, thick, AccentRed, dashed] (3,2.6) -- (3,1.4);
\draw[->, >=Stealth, thick, AccentRed, dashed] (9,3.8) -- (9,1.4);

\end{tikzpicture}
\caption{سرزمین‌های ازدست‌رفته‌ی ایران در عهدنامه‌های گلستان و ترکمنچای}
\label{fig:ch06-lost-territories}
\end{figure}

\subsection{نقش نجات‌طلبی مسیحیان قفقاز در از دست‌رفتن سرزمین‌ها}
\label{subsec:ch06-christian-role}

\begin{tcolorbox}[defbox, title={پرسش حساس: آیا نجات‌طلبی مسیحیان قفقاز باعث از دست‌رفتن قفقاز شد؟}]
پاسخ ساده نیست و باید چند لایه را از هم تفکیک کرد:

\begin{enumerate}[nosep, label=\textbf{\arabic*.}]
  \item \textbf{عامل ساختاری:} ضعف نظامی و فناوری ایران قاجار در مقابل روسیه‌ی مدرن‌شده — حتی بدون نجات‌طلبی مسیحیان، شکست نظامی محتمل بود.

  \item \textbf{عامل دموگرافیک:} جمعیت مسلمان قفقاز (آذری‌ها، لزگی‌ها، چچنی‌ها) اغلب در برابر روسیه مقاومت کردند — بنابراین نجات‌طلبی مختص مسیحیان بود و نه عمومی.

  \item \textbf{عامل سیاست ایران:} سیاست‌های تبعیض‌آمیز حکام محلی ایرانی نسبت به مسیحیان (مالیات‌های نابرابر، محدودیت‌های مذهبی)، زمینه‌ساز نجات‌طلبی بود.%
  \footnote{\lr{Bournoutian, George A. \textit{A Concise History of the Armenian People}. 6th ed. Mazda, 2012, pp.~225--248.}}

  \item \textbf{عامل ژئوپولیتیک:} روسیه در هر صورت به سمت قفقاز و آب‌های گرم حرکت می‌کرد — مسیحیان ابزار بودند نه علت.

  \item \textbf{نتیجه‌گیری بی‌طرفانه:} نجات‌طلبی مسیحیان \textbf{عامل تسهیل‌کننده} بود نه \textbf{عامل اصلی}. عامل اصلی، ضعف ساختاری ایران و توسعه‌طلبی روسیه بود. با این حال، نجات‌طلبی مسیحیان «مشروعیت» لازم را برای مداخله‌ی روسیه فراهم کرد.
\end{enumerate}
\end{tcolorbox}

% ─────────────────────────────────────────────────────────────────────────
\section{جنوب ایران: خلیج فارس، شیخ خزعل و بریتانیا}
\label{sec:ch06-persian-gulf}
% ─────────────────────────────────────────────────────────────────────────

\subsection{ساختار قدرت در خلیج فارس}
\label{subsec:ch06-gulf-structure}

خلیج فارس در دوره‌ی قاجار منطقه‌ای بود با:
\begin{itemize}[nosep, rightmargin=1em]
  \item حکومت مرکزی ضعیف (دسترسی محدود تهران به جنوب)
  \item شیوخ و خوانین محلی قدرتمند (بنی‌کعب، مرازیق، قواسم)
  \item حضور فزاینده‌ی بریتانیا (شرکت هند شرقی، سپس دولت بریتانیا)
  \item کشف نفت (۱۹۰۸) که اهمیت استراتژیک منطقه را دگرگون کرد
\end{itemize}

\subsection{شیخ خزعل: نجات‌طلبی یا استقلال‌طلبی؟}
\label{subsec:ch06-khazal}

\histref{شیخ خزعل خان کعبی}{ح. ۱۲۸۰–۱۳۱۵ ش / ۱۸۶۱–۱۹۳۶ م}، حاکم عربستان (خوزستان)، مهم‌ترین نمونه‌ی نجات‌طلبی/استقلال‌طلبی در جنوب ایران است. او با بهره‌گیری از ضعف حکومت مرکزی قاجار و سپس مخالفت با تمرکزگرایی رضاشاه، روابط گسترده‌ای با بریتانیا برقرار کرد.%
\footnote{کسروی، احمد. \textit{تاریخ پانصدساله‌ی خوزستان}. امیرکبیر، ۱۳۳۳، ص ۲۲۰--۲۸۰.}

\begin{longtable}{
  >{\raggedleft\arraybackslash\bfseries}p{2.5cm}
  >{\raggedleft\arraybackslash}p{5cm}
  >{\raggedleft\arraybackslash}p{5cm}}
\caption{مراحل رابطه‌ی شیخ خزعل با بریتانیا}\label{tab:ch06-khazal}\\
\toprule
\rowcolor{PrimaryDeep}
\textcolor{white}{\textbf{دوره}} &
\textcolor{white}{\textbf{ماهیت رابطه}} &
\textcolor{white}{\textbf{منفعت طرفین}} \\
\midrule
\endfirsthead
\toprule
\rowcolor{PrimaryDeep}
\textcolor{white}{\textbf{دوره}} &
\textcolor{white}{\textbf{ماهیت رابطه}} &
\textcolor{white}{\textbf{منفعت طرفین}} \\
\midrule
\endhead
\bottomrule
\endlastfoot

۱۸۹۷–۱۹۰۸ &
حمایت متقابل: خزعل امنیت شرکت نفت را تأمین می‌کرد و بریتانیا از استقلال عمل او حمایت می‌کرد. &
خزعل: تحکیم قدرت محلی. بریتانیا: دسترسی به نفت. \\
\rowcolor{NeutralLight}
۱۹۰۸–۱۹۱۸ &
اوج رابطه: خزعل عنوان \lr{KCIE} بریتانیایی دریافت کرد. امتیاز نفت عملاً در قلمرو او بود. &
خزعل: حمایت نظامی-سیاسی. بریتانیا: امنیت نفت.%
\footnote{\lr{Olson, William J. "The Genesis of the Anglo-Persian Agreement of 1919." In \textit{Rebirth of a Nation}, ed. Keddie. Yale UP, 1981.}} \\
۱۹۱۹–۱۹۲۴ &
تنش فزاینده: خزعل از تمرکزگرایی رضاخان نگران شد و از بریتانیا خواست مانع شود. بریتانیا «نجات‌طلبی» او را نشنیده گرفت. &
بریتانیا: رضاخان قوی‌تر و «مطمئن‌تر» از خزعل بود.%
\footnote{\lr{Cronin, Stephanie. \textit{The Army and the Creation of the Pahlavi State in Iran, 1910--1926}. I.B. Tauris, 1997, pp.~134--158.}} \\
\rowcolor{NeutralLight}
۱۹۲۵ &
سقوط: رضاخان خوزستان را فتح کرد. خزعل دستگیر و تبعید شد. بریتانیا مداخله نکرد. &
بریتانیا: منافع نفتی با رضاخان هم تأمین می‌شد. \\

\end{longtable}

\begin{tcolorbox}[enemybox, title={درس تاریخی: بی‌وفایی حامی خارجی}]
سرنوشت شیخ خزعل \textbf{نمونه‌ی کلاسیک بی‌وفایی حامی خارجی} است:
\begin{itemize}[nosep, rightmargin=1em]
  \item بریتانیا تا زمانی که خزعل مفید بود، از او حمایت کرد.
  \item هنگامی که بازیگر قدرتمندتری (رضاخان) ظهور کرد، بریتانیا خزعل را رها کرد.
  \item خزعل در تبعید در تهران در شرایط مشکوکی درگذشت (۱۳۱۵ ش / ۱۹۳۶ م).%
  \footnote{بیات، کاوه. «پایان کار شیخ خزعل.» \textit{گنجینه‌ی اسناد}، ش ۷۳--۷۴، ۱۳۸۸، ص ۵--۲۸.}
\end{itemize}

این الگو — \concept{حامی خارجی متحد محلی خود را فدای منافعش می‌کند} — در سراسر تاریخ تکرار شده: از لخمیان ساسانی تا مدعیان اشکانی در روم، و از خزعل تا بارزانی‌ها در کردستان عراق.
\end{tcolorbox}

% ─────────────────────────────────────────────────────────────────────────
\section{آذربایجان و قفقاز جنوبی: دست‌به‌دامان روسیه شدن؟}
\label{sec:ch06-azerbaijan}
% ─────────────────────────────────────────────────────────────────────────

\subsection{خانات آذربایجان: میان ایران و روسیه}
\label{subsec:ch06-khanates}

خانات آذربایجان (قفقاز جنوبی) در اواخر دوره‌ی قاجار، واحدهای نیمه‌مستقلی بودند که وفاداری‌شان میان ایران و روسیه در نوسان بود. \thinker{تادئوش سویه‌توخوفسکی} (\lr{Tadeusz Swietochowski}) نشان داده که:

\begin{itemize}[nosep, rightmargin=1em]
  \item \textbf{خانات مسلمان (شیعه):} مانند باکو، شیروان و لنکران — وفاداری‌شان به ایران قوی‌تر بود اما نه مطلق. برخی خان‌ها برای تحکیم قدرت محلی از هر دو طرف بهره می‌بردند.
  \item \textbf{خانات مسیحی:} مانند قره‌باغ (با جمعیت مختلط ارمنی-آذری) — ملیک‌های ارمنی اغلب از روسیه حمایت می‌کردند.%
  \footnote{\lr{Swietochowski, 1995, pp.~12--28.}}
  \item \textbf{الگوی عمومی:} خان‌ها «نجات‌طلب» نبودند بلکه \concept{فرصت‌طلب} بودند — به هر طرفی که در آن لحظه قوی‌تر بود تمایل نشان می‌دادند.
\end{itemize}

\begin{table}[htbp]
\centering
\caption{موضع خانات آذربایجان در جنگ اول ایران و روسیه}\label{tab:ch06-khanates}
\begin{adjustbox}{max width=\linewidth}
\begin{tabular}{
  >{\raggedleft\arraybackslash\bfseries}p{2.5cm}
  >{\raggedleft\arraybackslash}p{2cm}
  >{\raggedleft\arraybackslash}p{2.5cm}
  >{\raggedleft\arraybackslash}p{2.5cm}
  >{\raggedleft\arraybackslash}p{3cm}}
\toprule
\rowcolor{PrimaryDeep}
\textcolor{white}{خانات} &
\textcolor{white}{جمعیت غالب} &
\textcolor{white}{موضع اولیه} &
\textcolor{white}{موضع نهایی} &
\textcolor{white}{تحلیل} \\
\midrule

گنجه &
آذری (شیعه) &
مقاومت شدید &
سقوط (۱۸۰۴) &
وفاداری به ایران%
\footnote{\lr{Atkin, 1980, pp.~104--108.}} \\
\rowcolor{NeutralLight}
باکو &
آذری (شیعه) &
نوسان &
تسلیم (۱۸۰۶) &
فرصت‌طلبی \\
قره‌باغ &
مختلط &
موضع دوگانه &
تسلیم (۱۸۰۵) &
ملیک‌های ارمنی نقش تعیین‌کننده%
\footnote{\lr{Bournoutian, 1992, pp.~62--85.}} \\
\rowcolor{NeutralLight}
شکی &
آذری (سنّی) &
تسلیم سریع &
الحاق &
ضعف نظامی \\
تالش &
تالشی (شیعه) &
مقاومت &
تقسیم &
وفاداری نسبی به ایران \\
\rowcolor{NeutralLight}
ایروان &
مختلط (ارمنی-آذری) &
مقاومت ایرانی &
سقوط (۱۸۲۷) &
ارمنی‌ها از روسیه استقبال کردند \\

\bottomrule
\end{tabular}
\end{adjustbox}
\end{table}

% ─────────────────────────────────────────────────────────────────────────
\section{الگوهای استخراج‌شده از دوره‌ی صفوی-قاجار}
\label{sec:ch06-patterns}
% ─────────────────────────────────────────────────────────────────────────

\begin{tcolorbox}[wavebox, title={الگوهای کلیدی دوره‌ی صفوی-قاجار}]
\begin{enumerate}[nosep, label=\textbf{الگوی \arabic*:}]
  \item \textbf{نجات‌طلبی مذهبی:} هم‌کیشی (مسیحیت) مهم‌ترین عامل نزدیکی مرزنشینان قفقازی به روسیه بود — نه نارضایتی اقتصادی یا سیاسی صرف. \\[4pt]

  \item \textbf{ضعف مرکز = افزایش نجات‌طلبی:} هر بار که حکومت مرکزی ایران ضعیف شد (فروپاشی صفویه، هرج‌ومرج افشاری-زندی، ضعف قاجار)، مرزنشینان بیشتر به بیگانه متوسل شدند. \\[4pt]

  \item \textbf{«بیگانه‌ی قابل‌دسترس»:} ظهور روسیه‌ی تزاری به‌عنوان قدرت مسیحی همسایه، برای نخستین بار «بیگانه‌ی قابل‌دسترسی» در اختیار مسیحیان قفقاز قرار داد — در دوره‌ی پیش از صفوی، چنین گزینه‌ای وجود نداشت. \\[4pt]

  \item \textbf{بی‌وفایی حامی خارجی:} در تمامی موارد بررسی‌شده، حامی خارجی (روسیه، بریتانیا) منافع خود را بر خواسته‌ی دعوت‌کنندگان ترجیح داد — عهدنامه‌ی گئورگیفسک (وعده‌ی حمایت بدون عمل) و شیخ خزعل (رهاشدن در بزنگاه). \\[4pt]

  \item \textbf{نتیجه‌ی نهایی: الحاق نه استقلال:} گرجستان به‌جای «متحد آزاد روسیه»، استان روسیه شد. ارمنستان شرقی به‌جای «کشور مستقل»، گوبرنیای روسیه شد. خزعل به‌جای «حاکم مستقل»، زندانی رضاشاه شد. \\[4pt]

  \item \textbf{عدم تقارن اطلاعاتی:} دعوت‌کنندگان اغلب از منافع واقعی مداخله‌گر آگاه نبودند (یا نمی‌خواستند آگاه باشند) — نمونه‌ی بارز \concept{آرزواندیشی}.
\end{enumerate}
\end{tcolorbox}

% ─────────────────────────────────────────────────────────────────────────
\section{نمودار ژئوپولیتیک: ایران در محاصره‌ی قدرت‌های بزرگ (۱۸۰۰–۱۹۲۵)}
\label{sec:ch06-geopolitical}
% ─────────────────────────────────────────────────────────────────────────

\begin{figure}[htbp]
\centering
\begin{tikzpicture}[
  every node/.style={font=\small},
  empire/.style={
    draw=PrimaryMid, fill=BgCool!30, rounded corners=6pt,
    text width=3cm, align=center, minimum height=1.5cm,
    font=\small\bfseries
  },
  iran/.style={
    draw=AccentGold, fill=AccentGold!15, rounded corners=8pt,
    text width=3.5cm, align=center, minimum height=2cm,
    font=\normalsize\bfseries
  },
  buffer/.style={
    draw=NeutralMid, fill=NeutralLight, rounded corners=3pt,
    text width=2.5cm, align=center, minimum height=0.9cm,
    font=\scriptsize
  },
  influence/.style={->, >=Stealth, thick, AccentRed},
  resist/.style={->, >=Stealth, thick, AccentGreen, dashed},
]

% ایران مرکز
\node[iran] (iran) at (7,4) {\rl{ایران قاجار}\\%
  \rl{\scriptsize ضعیف و در حال تجزیه}};

% قدرت‌ها
\node[empire, fill=blue!10] (russia) at (7,8.5) {\rl{روسیه‌ی تزاری}};
\node[empire, fill=red!10] (britain) at (13,4) {\rl{بریتانیا}};
\node[empire, fill=green!10] (ottoman) at (1,4) {\rl{عثمانی}};

% مناطق بافر / نجات‌طلبی
\node[buffer] (georgia) at (3,7) {\rl{گرجستان}};
\node[buffer] (armenia) at (5,7) {\rl{ارمنستان شرقی}};
\node[buffer] (azer) at (9,7) {\rl{خانات آذربایجان}};
\node[buffer] (khuz) at (11,2) {\rl{خوزستان (خزعل)}};
\node[buffer] (baluch) at (11,6) {\rl{بلوچستان}};
\node[buffer] (kurd) at (2,2) {\rl{کردستان}};
\node[buffer] (herat) at (12,7.5) {\rl{هرات و افغانستان}};

% فلش‌های نفوذ
\draw[influence] (russia) -- (georgia);
\draw[influence] (russia) -- (armenia);
\draw[influence] (russia) -- (azer);
\draw[influence] (britain) -- (khuz);
\draw[influence] (britain) -- (baluch);
\draw[influence] (britain) -- (herat);
\draw[influence] (ottoman) -- (kurd);

% فلش‌های مقاومت
\draw[resist] (iran) -- (georgia);
\draw[resist] (iran) -- (armenia);
\draw[resist] (iran) -- (azer);
\draw[resist] (iran) -- (khuz);
\draw[resist] (iran) -- (kurd);

\end{tikzpicture}
\caption{ایران در محاصره‌ی قدرت‌ها و مناطق نجات‌طلبی (۱۸۰۰–۱۹۲۵)}
\label{fig:ch06-geopolitical}
\end{figure}

% ─────────────────────────────────────────────────────────────────────────
\section{مقایسه‌ی تطبیقی: قفقاز و خلیج فارس}
\label{sec:ch06-comparison}
% ─────────────────────────────────────────────────────────────────────────

\begin{table}[htbp]
\centering
\caption{مقایسه‌ی نجات‌طلبی در قفقاز و خلیج فارس}\label{tab:ch06-comparison}
\begin{adjustbox}{max width=\linewidth}
\begin{tabular}{
  >{\raggedleft\arraybackslash\bfseries}p{3cm}
  >{\raggedleft\arraybackslash}p{5cm}
  >{\raggedleft\arraybackslash}p{5cm}}
\toprule
\rowcolor{PrimaryDeep}
\textcolor{white}{ویژگی} &
\textcolor{white}{قفقاز (گرجی‌ها/ارمنی‌ها)} &
\textcolor{white}{خلیج فارس (شیخ خزعل)} \\
\midrule

عامل اصلی نجات‌طلبی &
هویت مذهبی (مسیحیت) + فشار نظامی &
قدرت‌طلبی شخصی + ضعف مرکز \\
\rowcolor{NeutralLight}
حامی خارجی &
روسیه‌ی تزاری &
بریتانیا \\
نوع رابطه &
ایدئولوژیک-مذهبی + نظامی &
اقتصادی-عملی (نفت) \\
\rowcolor{NeutralLight}
پشتیبانی مردمی &
نسبتاً گسترده (جمعیت مسیحی) &
محدود (قبایل عرب بنی‌کعب) \\
نتیجه‌ی نهایی &
الحاق به روسیه &
سرکوب توسط رضاشاه \\
\rowcolor{NeutralLight}
پیامد بلندمدت &
استقلال ۱۹۱۸ (موقت)، شوروی، ۱۹۹۱ (دائم) &
ادغام کامل در ایران \\
درجه‌ی «موفقیت» &
D (الحاق، نه استقلال) &
F (شکست کامل) \\

\bottomrule
\end{tabular}
\end{adjustbox}
\end{table}

% ─────────────────────────────────────────────────────────────────────────
\section{مشروطه و مداخله‌ی خارجی: الگوی نوین}
\label{sec:ch06-constitutional}
% ─────────────────────────────────────────────────────────────────────────

\subsection{اولتیماتوم روسیه و بمباران مجلس (۱۲۹۰ ش / ۱۹۱۱ م)}
\label{subsec:ch06-ultimatum}

دوره‌ی مشروطه (۱۲۸۵–۱۲۹۰ ش / ۱۹۰۶–۱۹۱۱ م) الگوی جدیدی از تعامل با قدرت خارجی ارائه داد: مشروطه‌خواهان از یک‌سو به اصول حقوق بین‌المللی و حمایت افکار عمومی جهانی متوسل شدند، و از سوی دیگر با مداخله‌ی مستقیم روسیه و بریتانیا مواجه بودند.%
\footnote{آفاری، ژانت. \textit{انقلاب مشروطه‌ی ایران}. ترجمه‌ی رضا رضایی. بیستون، ۱۳۸۵.}

اولتیماتوم روسیه در آبان ۱۲۹۰ ش / نوامبر ۱۹۱۱ م — که خواستار اخراج \thinker{مورگان شوستر} (\lr{Morgan Shuster})، مستشار مالی امریکایی شد — نمونه‌ی بارز مداخله‌ی خارجی «بدون دعوت» بود. نکته‌ی تلخ آنکه \textbf{بریتانیا — که مشروطه‌خواهان امید به حمایتش داشتند — با روسیه هم‌صدا شد.}%
\footnote{\lr{Shuster, W. Morgan. \textit{The Strangling of Persia}. Century Co., 1912. Reprinted: Mage, 2006.}}

\begin{histQuote}{مورگان شوستر، ۱۹۱۲}
«خفه‌کردن ایران بهترین توصیف از آن چیزی است که بریتانیا و روسیه با رضایت متقابل با این کشور کردند ... ایرانیان فریاد دادخواهی سر دادند اما هیچ‌کس نشنید.»%
\footnote{\lr{Shuster, 1912/2006, p.~xxii.}}
\end{histQuote}

\subsection{قرارداد ۱۹۰۷ و تقسیم ایران}
\label{subsec:ch06-1907}

\begin{tcolorbox}[enemybox, title={قرارداد ۱۹۰۷: تقسیم ایران بدون نجات‌طلبی}]
قرارداد ۱۹۰۷ انگلیس و روسیه، ایران را بدون رضایت یا درخواست هیچ‌بخشی از مردم ایران، به سه منطقه‌ی نفوذ تقسیم کرد:
\begin{itemize}[nosep, rightmargin=1em]
  \item \textbf{شمال:} منطقه‌ی نفوذ روسیه (شامل تهران، تبریز، اصفهان)
  \item \textbf{جنوب‌شرق:} منطقه‌ی نفوذ بریتانیا (شامل سیستان، بلوچستان)
  \item \textbf{مرکز:} منطقه‌ی «بی‌طرف» (که عملاً بازار رقابت دو قدرت بود)
\end{itemize}

این قرارداد نشان‌دهنده‌ی آن است که مداخله‌ی خارجی \textbf{لزوماً نیازمند «دعوت» از داخل نیست} — قدرت‌های بزرگ می‌توانند بدون هیچ درخواست داخلی، سرنوشت ملتی را تعیین کنند.%
\footnote{\lr{Kazemzadeh, Firuz. \textit{Russia and Britain in Persia, 1864--1914}. Yale UP, 1968, pp.~458--502.}}
\end{tcolorbox}

% ─────────────────────────────────────────────────────────────────────────
\section{جمع‌بندی فصل}
\label{sec:ch06-summary}
% ─────────────────────────────────────────────────────────────────────────

\begin{tcolorbox}[wavebox, title={یافته‌های کلیدی فصل ۶}]
\begin{enumerate}[nosep, label=\textbf{\arabic*.}]
  \item دوره‌ی صفوی-قاجار \textbf{نخستین دوره‌ی مستند نجات‌طلبی سازمان‌یافته} در تاریخ ایران است — بر خلاف دوران باستان که بیشتر «تمکین عملی» بود.

  \item \textbf{هم‌کیشی مذهبی} (مسیحیت ارتدوکس) مهم‌ترین عامل تسهیل‌کننده‌ی نجات‌طلبی مرزنشینان قفقازی بود. در غیاب این عامل (مثلاً در مورد آذری‌های شیعه)، نجات‌طلبی بسیار کم‌تر بود.

  \item \textbf{هیچ‌یک از نجات‌طلبی‌های بررسی‌شده به استقلال واقعی نینجامید:}
  \begin{itemize}[nosep]
    \item گرجستان: الحاق به روسیه (۱۸۰۱)
    \item ارمنستان شرقی: گوبرنیای روسیه (۱۸۲۸)
    \item خزعل: سرکوب و تبعید (۱۹۲۵)
  \end{itemize}

  \item \textbf{الگوی «بی‌وفایی حامی خارجی»} بار دیگر تأیید شد: روسیه (گرجستان)، بریتانیا (خزعل) و هر دو (مشروطه).

  \item \textbf{قرارداد ۱۹۰۷} نشان داد که مداخله‌ی خارجی لزوماً نیازمند «دعوت» نیست — امپریالیسم مستقل از نجات‌طلبی عمل می‌کند.

  \item \textbf{درس اصلی:} ضعف حکومت مرکزی، بزرگ‌ترین عامل آسیب‌پذیری در برابر هم نجات‌طلبی مرزنشینان و هم مداخله‌ی یک‌جانبه‌ی قدرت‌های بزرگ است.
\end{enumerate}
\end{tcolorbox}

\cleardoublepage